
\needspace{4\baselineskip}
\color{uniblue}{\subsection{Автоматическая частотная разгрузка (АЧР)}}
\color{black}

\begin{enumerate}

\item
АЧР предназначена для отключения нагрузки при снижении частоты в энергосистеме. В устройстве автоматическая частотная разгрузка представлена функцией «АЧР».

\item
Основные характеристики функции:
\begin{itemize}
    \item погрешность ИО частоты не более 20 мГц;
    \item погрешность срабатывания ИО скорости изменения частоты не более 150 мГц;
    \item погрешность ИО частоты в рабочем диапазоне при напряжении входного сигнала 100 В не более 10 мГц; 
    \item время срабатывания и возврата ИО частоты не более 60 мс;
    \item время срабатывания ИО скорости снижения частоты не более 200 мс.
\end{itemize}

\item
В функции предусмотрена настройка, позволяющая работать как от внешней АЧР, так и от своих измерительных органов. Выбор типа контролируемой АЧР осуществляется уставкой SGF2.

\item
Если выбран режим «По внешнему сигналу», то при наличии входного дискретного сигнала «АЧР» и отсутствии дискретного сигнала «Блокировка АЧР» формируются пусковые условия функции и начинается
отсчет выдержки времени таймера Т1 по истечении которого, формируется сигнал «АЧР: Отключение».

\item
Если выбран режим «Измеренный», то функция реагирует на снижение максимального значения частот ниже уставки и формирует сигнал пуска, который блокируется по одному из двух условий:
\begin{itemize}
    \item снижение линейного напряжения ниже уставки;
    \item снижение скорости изменения частоты ниже уставки, если уставка SGF3 введена в работу.
\end{itemize}

\item
Для отстройки от дребезга порогового элемента и продления сигнала от ИО частоты в функции предусмотрена регулируемая выдержка времени на возврат, задаваемая уставкой T2.

\item
В функции предусмотрен выходной сигнал «АЧР: Отключить», работающий в импульсном режиме длительностью таймера Т3, который воздействует непосредственно на выходные реле для отключения выключателя.

\item
Возврат АЧР происходит при превышении частоты хотя бы в одной из трех фаз выше значения задаваемого уставкой.

\item
Ввод функции в работу на этапе параметрирования устройства осуществляется программным переключателем SGF1 «Ввод\_функции». Информация о введенной функции отображается выходным сигналом «АЧР: Ввод».

\item
Функция может быть выведена из работы оперативно путем активации сигнала «ОВ АЧР», что характеризуется активным состоянием сигнала «АЧР: Оперативный вывод» на выходе функции. 

\item
Алгоритм работы функции выполнен в соответствии с рисунком \ref{achr:img1}. В таблице \ref{achr:tbl1} приведены параметры, необходимые для настройки функции.

\vspace{1mm}
\begin{figure}[H]
\centering
\includegraphics[width=1\textwidth,height=1\textheight,keepaspectratio]{\fbpath/119. АЧР/_latex/img/achr.pdf}
\captionof{figure}{Функциональная схема логики АЧР}\label{achr:img1}
\end{figure}

\end{enumerate}
